\section{Results}

In this section, I will present the results of the experiments previous conduced.
The results are summarized in Table \ref{tab:results}. 
The table includes the mean and standard deviation (SD) of the F1-Score, ROC AUC, and AUPRC for each method evaluated. 
The methods compared include MCD, k-NN (unsupervised), k-NN (semi-supervised), and InterCont.
After presenting the results, I will display the confusion matrices for each method.
Additionally I will plot the the AUPRC and ROC AUC to further analyse the performance of each method.


\begin{table}[h]
\centering
\caption{Experimental Results}
\label{tab:results}
\begin{tabular}{lcccccc} % l=left aligned, r=right aligned
\toprule
\textbf{Method}         & \textbf{F1-Score} & \textbf{SD} & \textbf{ROC AUC} & \textbf{SD} & \textbf{AUPRC} & \textbf{SD}          \\
\midrule
MCD                     & 50.00             &       & \textbf{97.56}    &          & 71.04             &        \\
k-NN (unsupervised)     & 50.00             &       &         94.29     &          & 49.83             &        \\
k-NN (semi-supervised)  & \textbf{89.07}    & 3.97  &         96.12     & 2.42     & 88.03             & 8.36   \\
InterCont               & 77.44             & 10.59 &         96.70     & 2.06     & \textbf{99.48}    & 0.32   \\
\bottomrule
\end{tabular}
\end{table}


The best performances are written in bold letters. 
The method with the highest F1-Score is the semi-supervised k-NN method with 89.07\%, while the method with the highest ROC AUC is MCD with 97.56\%.
InterCont performs best in terms of AUPRC, with a resulting score of 99.48\%.
The differences between the methods vary alot.
The difference of the best to worst performance in terms of F1-Score, 79.07 percentage points.
AUPRC differs by 69.46 percentage points between the best and worst performing method.
While the difference between the best and worst performance in terms of ROC AUC is only 13.20 percentage points.


The standard deviations differ a lot between the methods as well.
MCD has no standard deviation, as it is a deterministic method.
The unsupervised k-NN method has the highest deviation in terms of F1-Score, with 20 percentage points, while the semi-supervised k-NN method has the lowest deviation with 3.97 percentage points.
InterCont deviates by 10.59 percentage points in terms of F1-Score.
The standard deviation of ROC AUC are not differing so much between methods as F1-Score or AUPRC.
The unsupervised k-NN method has the highest deviation with 4.83 percentage points, while InterCont has the lowest deviation with 2.06 percentage points.
The semi-supervised k-NN method has a deviation of 2.42 percentage points, which is only slightly higher than InterCont.
The standard deviation of AUPRC is highest for the unsupervised k-NN method with 12.49 percentage points, while InterCont has the lowest deviation with 0.32 percentage points.
The semi-supervised k-NN method has a deviation of 8.36 percentage points, which is higher than InterCont but much lower than the unsupervised k-NN method.


The authors of InterCont tested their results with a wilcoxon signed rank test. 
They conducted that by statistical sinificance, that their model outperforms other methods.
I can't apply a wilcoxon signed rank test to the results of my experiments, because this I applied only one dataset.
Therefore the num






\section{Discussion}

As a general detection performance result, we can conclude that nearest-neighbor based
algorithms perform better in most cases when compared to clustering algorithms. Also, the sta-
bility concerning a not-perfect choice of k is much higher for the nearest-neighbor based meth-
ods. The reason for the higher variance in clustering-based algorithms is very likely due to the
non-deterministic nature of the underlying k-means clustering algorithm. Despite of this
disadvantage, clustering-based algorithms have a lower computation time. As a conclusion, we
recommend to prefer nearest-neighbor based algorithms if computation time is not an issue. If
a faster computation is required for large datasets, for example in a near real-time setting, clus-
tering-based anomaly detection might be the method of choice. For small datasets, clustering-
based methods should be avoided.Among the nearest-neighbor based methods, the global k-NN algorithm is a good candidate
on average.


\todo{insert analysis of the outliers}
\todo{explain why MCD is bad for real world data, why k-NN is better for multimodel clusters. Why is contrastive learning the best choice for the real world dataset}
\todo{the resulting intercont scores are interpretable, k-NN is the distance interpretable}
